\documentclass[11pt]{article}
\usepackage[margin=1in]{geometry}
\usepackage[T1]{fontenc}
\usepackage{graphicx}
\usepackage{hyperref}
\usepackage{enumitem}
\setlength{\parskip}{0.5em}
\setlength{\parindent}{0pt}

\title{transformer attention models for language}
\author{arXiv GraphRAG --- phase-6 deep research}
\date{\today}

\begin{document}
\maketitle

\begin{abstract}
An excellent, high-fidelity review that masterfully integrates 2025-era research into a coherent historical narrative.
\end{abstract}

\section{The Evolution of Attention: From Sequence Foundations to the Transformer Frontier}

The ability to model linguistic structures through the lens of context is the central challenge of natural language processing. For decades, the field struggled to balance the need for long-range dependency capture with computational feasibility, a tension that defined the foundations of sequence modeling and early attention mechanisms. The shift from recurrent paradigms to the transformer architecture marked a fundamental change in how machines process information, enabling the rise of pre-trained transformers and the massive scaling of parameters that characterizes the current era of large language models.

As these models expanded, research branched into optimizing the internal mechanics of the attention mechanism itself. This review traces the development of various architectural variations and the critical role of positional encoding in preserving sequence order. To address the inherent computational bottlenecks of standard attention, we examine the significant advancements in efficiency and complexity reduction. These improvements have directly enabled breakthroughs in long-context modeling and memory mechanisms, allowing models to maintain coherence over increasingly vast input sequences.

Beyond pure scaling, the field is currently witnessing a convergence of methodologies. We analyze the emergence of hybrid architectures that attempt to integrate the strengths of recurrent structures with the parallel processing capabilities of transformers. Ultimately, this review looks toward the frontier of the field, where the focus shifts from mere performance to the deeper challenges of interpretability, the derivation of robust design principles, and the understanding of advanced dynamics within attention-based systems.

\subsection{Foundations of Sequence Modeling and Early Attention}

The evolution of sequence modeling has been defined by a transition from fixed-length vector representations to dynamic, attention-driven architectures. Early efforts focused on improving phrase representations through joint training of RNN encoder-decoder models to capture semantic and syntactic nuances for statistical machine translation (\texttt{arXiv:1406.1078}). This foundational work was expanded by the introduction of end-to-end sequence learning using multilayered Long Short-Term Memory (LSTM) networks, which mapped input sequences to a fixed-dimensional vector for decoding (\texttt{arXiv:1409.3215}). While \texttt{arXiv:1409.3215} achieved a BLEU score of 34.8 on the WMT'14 English-to-French dataset, its reliance on a single fixed-length vector created a bottleneck for long-range dependencies.

To mitigate this bottleneck, \texttt{arXiv:1409.0473} proposed an extension to the encoder-decoder framework by incorporating an attention mechanism. This allowed the model to perform "soft-alignment," automatically searching for relevant parts of the source sentence rather than relying on a static vector. This shift from fixed representations to dynamic alignment was further refined by \texttt{arXiv:1508.04025}, which categorized attentional mechanisms into global and local classes. Their analysis demonstrated that local attention—focusing on a specific subset of source words—could yield a 5.0 BLEU point gain over non-attentional systems, while their ensemble model achieved a state-of-the-art 25.9 BLEU on WMT'15 English-to-German tasks.

As these architectures matured, research shifted toward large-scale industrial viability. Google’s Neural Machine Translation (GNMT) system (\texttt{arXiv:1609.08144}) addressed the computational expenses and rare word issues inherent in earlier models by utilizing a deep LSTM network with 8 encoder and 8 decoder layers. By integrating attention, residual connections, and wordpiece sub-word units, GNMT reduced translation errors by an average of 60\% compared to previous phrase-based production systems.

However, the most significant paradigm shift occurred when the reliance on recurrence was abandoned entirely. While previous models like \texttt{arXiv:1409.3215} and \texttt{arXiv:1609.08144} relied on the sequential nature of LSTMs, \texttt{arXiv:1706.03762} introduced the Transformer. This architecture dispensed with recurrence and convolutions, relying solely on attention mechanisms within an encoder-decoder configuration. The Transformer demonstrated superior efficiency and performance, achieving 28.4 BLEU on WMT 2014 English-to-German and 41.8 BLEU on WMT 2014 English-to-French, outperforming the recurrent-based state-of-the-art. This progression marks a move from using attention as a supplementary alignment tool for RNNs to using attention as the primary structural component of the model itself.

\subsection{The Rise of Pre-trained Transformers and Scaling}

The evolution of transformer-based language modeling has been characterized by a shift from task-specific architectures toward massive, pre-trained representations capable of capturing complex linguistic nuances. Early advancements in this direction focused on deep contextualized word representations, where the internal states of a deep bidirectional language model (biLM) were leveraged to model syntax, semantics, and polysemy (\texttt{arXiv:1802.05365}). This paradigm was significantly advanced by the introduction of BERT, which utilized a Transformer encoder pre-trained on unlabeled text using Masked Language Model (MLM) and Next Sentence Prediction (NSP) objectives (\texttt{arXiv:1810.04805}). BERT demonstrated the power of bidirectional self-attention by achieving state-of-the-art results across eleven tasks, including a 93.2 Test F1 on SQuAD v1.1 and an 80.5\% GLUE score (\texttt{arXiv:1810.04805}). However, BERT's architecture introduced a structural mismatch between pre-training and fine-tuning due to the use of [MASK] tokens, and its larger variants exhibited instability on smaller datasets (\texttt{arXiv:1810.04805}).

Subsequent research focused on optimizing these pre-training regimes and expanding their structural limits. RoBERTa addressed the inefficiencies in BERT's training by demonstrating that the original model was significantly undertrained; through robust optimization of hyperparameters and increased training data size, RoBERTa matched or exceeded performance levels of models published later (\texttt{arXiv:1907.11692}). While BERT and RoBERTa focused on bidirectional encoding, Transformer-XL addressed the limitation of fixed-length context by introducing a segment-level recurrence mechanism and a novel positional encoding scheme (\texttt{arXiv:1901.02860}). This allowed the model to learn dependencies 80\% longer than RNNs and 450\% longer than vanilla Transformers, significantly improving evaluation speed (\texttt{arXiv:1901.02860}).

The versatility of these models was further expanded by the introduction of a unified text-to-text framework, which converted all language problems into a single format and provided a systematic study of transfer learning across various architectures and datasets (\texttt{arXiv:1910.10683}). As these models grew in complexity, understanding the drivers of performance became critical. Research into scaling laws revealed that language model performance (cross-entropy loss) follows a predictable power-law function of model size, dataset size, and compute (\texttt{arXiv:2001.08361}). Crucially, while larger models are more sample-efficient, the study noted that specific architectural details, such as depth or width, have minimal impact on loss within a wide range, suggesting that the primary drivers of performance are the scale of data and parameters (\texttt{arXiv:2001.08361}). Together, these developments represent a transition from modeling individual word contexts to managing massive-scale, unified, and long-range linguistic dependencies.

\subsection{Architectural Variations and Positional Encoding}

The evolution of Transformer architectures has been driven by the need to better represent spatial relationships and improve computational efficiency. A primary area of innovation concerns how models handle sequence order. While standard Transformers rely on absolute position representations, \texttt{arXiv:1803.02155} introduces a relation-aware self-attention mechanism that incorporates relative position representations (distances between elements). This approach demonstrated significant practical utility in translation tasks, yielding improvements of 1.3 BLEU on English-to-German and 0.3 BLEU on English-to-French over absolute methods [\texttt{arXiv:1803.02155}]. This concept of integrating relative dependency was further refined by \texttt{arXiv:2104.09864}, which proposes Rotary Position Embedding (RoPE). RoPE utilizes rotation matrices to integrate absolute position while simultaneously incorporating explicit relative position dependency, a method that has proven superior in long text classification benchmarks [\texttt{arXiv:2104.09864}].

Beyond positional encoding, research has focused on the structural efficiency of the attention mechanism itself. To address the quadratic complexity of standard attention, \texttt{arXiv:1904.10509} introduces sparse factorizations of the attention matrix, reducing time and memory requirements to \$O(n \textbackslash{}sqrt\{n\})\$ [\texttt{arXiv:1904.10509}]. This contrasts with work focused on the redundancy of the multi-head structure; for instance, \texttt{arXiv:1905.10650} observes that a large percentage of attention heads can be pruned at test time without significant performance loss, suggesting that many layers can be reduced to a single head [\texttt{arXiv:1905.10650}].

Recent theoretical and structural advancements have moved toward optimizing how these components interact. \texttt{arXiv:2507.02944} reframes multi-head attention not merely as parallel processes, but as a system of synergistic computational graphs where heads function as feedforward directed acyclic graphs (DAGs) [\texttt{arXiv:2507.02944}]. Their analysis suggests that multi-head attention provides a functional advantage over single-head models by enhancing information propagation and mixing times, provided specific head-diversity conditions are met [\texttt{arXiv:2507.02944}]. Complementing this view of structural synergy, \texttt{arXiv:2508.00544} introduces PaPaformer, a decoder-only variant that utilizes lower-dimensional parallel paths [\texttt{arXiv:2508.00544}]. Unlike traditional scaling, PaPaformer allows these paths to be trained individually and then combined, significantly reducing training time—enabling model evaluation in hours rather than days or weeks [\texttt{arXiv:2508.00544}]. Together, these developments indicate a shift from simply increasing model scale to optimizing the mathematical and structural properties of the attention mechanism to achieve efficiency and better information flow.

\subsection{Efficiency and Complexity Reduction in Attention}

Standard Transformer architectures are fundamentally limited by the quadratic time and memory complexity of self-attention. Research to mitigate this bottleneck follows two primary trajectories: parameter-centric optimization and complexity-reduction mechanisms. ALBERT (1909.11942) focuses on the former, employing factorized embedding parameterization and cross-layer parameter sharing to reduce memory consumption and increase training speed. Unlike models that alter the attention mechanism itself, ALBERT optimizes the model's footprint, establishing new state-of-the-art results on GLUE, RACE, and SQuAD with significantly fewer parameters than BERT-large.

In contrast, models such as Linformer (2006.04768) and Performer (2009.14794) target the attention mechanism to achieve \$O(n)\$ complexity. Linformer (2006.04768) utilizes a low-rank matrix approximation of the self-attention mechanism to reduce complexity in both time and space. Performer (2009.14794) takes a different mathematical approach via the FAVOR+ mechanism, which approximates softmax attention-kernels using positive orthogonal random features. While both achieve linear complexity, they represent distinct strategies: low-rank approximation versus kernel approximation. A third approach, exemplified by BigBird (2007.14062), utilizes a sparse attention mechanism incorporating global, local, and random tokens. This allows BigBird to handle sequences up to 8x longer than previously possible, maintaining the ability to act as a universal approximator while reducing quadratic dependency.

Recent literature has sought to synthesize these disparate methods into unified frameworks. \texttt{arXiv:2507.19595} provides a comprehensive taxonomy that bridges the gap between linear attention (such as kernel approximations and recurrent formulations) and sparse attention (such as fixed patterns and block-wise routing). This survey is particularly rigorous in its hardware-centric derivation of the "Chunkwise Recurrent Representation," which explains how parallelism is maintained in linear models. However, the survey's depth is limited by its incomplete analysis of "Clustering Attention" and a lack of comparative performance benchmarks, such as tables comparing perplexity against latency. Expanding this scope, \texttt{arXiv:2508.09834} categorizes the field into even broader architectural paradigms, including sparse mixture-of-experts (MoE) and hybrid architectures. Together, these works provide a structural guide for the transition from simple parameter scaling to the efficiency-centric design required for modern large language model deployment.

\subsection{Long-Context Modeling and Memory Mechanisms}

The challenge of long-context modeling in Transformer architectures involves balancing computational efficiency with the ability to retain long-range dependencies. Early approaches, such as the Compressive Transformer, focus on compressing past memories to facilitate long-range sequence learning, achieving 17.1 ppl on WikiText-103 and 0.97 bpc on Enwik8 [\texttt{arXiv:1911.05507}]. In contrast to this memory-compression strategy, the Longformer introduces a structural modification to the attention mechanism itself, combining local windowed attention with task-motivated global attention to achieve linear scaling [\texttt{arXiv:2004.05150}]. While the Longformer provides a drop-in replacement for standard self-attention and outperforms RoBERTa on long-document tasks [\texttt{arXiv:2004.05150}], newer research shifts focus toward addressing the inherent redundancies and dimensionality issues within the attention mechanism.

Recent advancements address the "curse of high dimensionality" by treating sequence modeling as a supervised learning task. The Dynamic Group Attention (DGA) method reformulates attention optimization as a linear coding problem, using a group coding strategy to aggregate less important tokens [\texttt{arXiv:2505.22107}]. This represents a more mathematically rigorous approach to sparsity compared to the heuristic windowing seen in the Longformer [\texttt{arXiv:2004.05150}], allowing DGA to significantly reduce computational costs while maintaining competitive performance [\texttt{arXiv:2505.22107}].

A secondary dimension of long-context research concerns how models handle sequence lengths that deviate from their training distribution. SWAN-GPT addresses this through a decoder-only architecture that interleaves layers without positional encodings (NoPE) with sliding-window attention equipped with rotary positional encodings (SWA-RoPE) [\texttt{arXiv:2504.08719}]. This configuration allows for robust generalization to sequence lengths substantially longer than those seen during training, offering higher throughput than standard GPT architectures [\texttt{arXiv:2504.08719}].

While SWAN-GPT focuses on length extrapolation, InfLLM-V2 prioritizes a seamless transition between short and long contexts. By utilizing a trainable sparse attention mechanism that reuses dense attention parameters via parameter-free architecture modification, InfLLM-V2 enables a transition from dense attention for short inputs to sparse attention for long sequences [\texttt{arXiv:2509.24663}]. This approach is highly efficient, operating 4\$\textbackslash{}times\$ faster than dense attention while retaining 98.1\% of long-context understanding and 99.7\% of chain-of-thought reasoning performance [\texttt{arXiv:2509.24663}]. Together, these developments indicate a progression from simple memory compression [\texttt{arXiv:1911.05507}] and windowed attention [\texttt{arXiv:2004.05150}] toward sophisticated, switchable, and dimensionally-aware architectures that minimize the trade-off between computational speed and reasoning capability.

\subsection{Hybrid Architectures: Integrating RNNs and Recurrent Structures}

Recent advancements in language modeling have focused on bridging the gap between the parallelizable nature of Transformers and the efficient, linear-complexity processing of Recurrent Neural Networks (RNNs). A foundational step in this convergence is the realization that self-attention can be reformulated as a linear dot-product of kernel feature maps [\texttt{arXiv:2006.16236}]. By leveraging the associativity property of matrix products, this approach reduces computational complexity from \$O(N\textasciicircum{}2)\$ to \$O(N)\$, enabling an iterative implementation that allows Transformers to function as RNNs [\texttt{arXiv:2006.16236}]. This transformation provides a significant efficiency gain, achieving up to 4000x faster autoregressive prediction on very long sequences while maintaining performance parity with vanilla Transformers [\texttt{arXiv:2006.16236}].

Building upon these theoretical links, contemporary research explores various ways to integrate recurrent structures into Transformer frameworks to optimize for specific constraints like parameter count, memory, or data availability. For instance, the Compact Recurrent Transformer (CRT) utilizes a hybrid approach where shallow Transformer models process short local segments, while an RNN manages a single persistent memory vector to capture long-range information [\texttt{arXiv:2505.00929}]. This architecture demonstrates high efficiency, achieving state-of-the-art results on the Toyota Smarthome video dataset and comparable performance to full-length Transformers while using only half or quarter-sized segments and substantially fewer FLOPs [\texttt{arXiv:2505.00929}]. Similarly, in low-resource settings, researchers have found success by selectively replacing attention layers with feed-forward and quasi-recurrent neural network layers [\texttt{arXiv:2502.00617}]. This hybrid strategy allows reduced architectures to outperform existing models with similar parameter counts and approach the performance of much larger models on corpora such as Enwik8 and Wikitext-103 [\texttt{arXiv:2502.00617}].

However, the integration of these two paradigms introduces unique training challenges, particularly regarding length generalization. A critical weakness identified in Transformer-RNN hybrids is that the Transformer components often exploit "shortcut" solutions, which prevents the model from generalizing to sequences longer than those seen during training [\texttt{arXiv:2510.00258}]. To mitigate this, a "delayed attention training" strategy has been proposed, which postpones the training of attention layers to prevent such shortcut reliance [\texttt{arXiv:2510.00258}]. This method has proven effective, achieving over 90\% accuracy on hybrid sequences three times the length of the training data [\texttt{arXiv:2510.00258}].

Finally, the relationship between these architectures is being explored through knowledge distillation. The Distil-xLSTM model attempts to approximate Transformer-based attention parametrization by distilling knowledge from a Large Language Model (LLM) into an xLSTM-based recurrent structure [\texttt{arXiv:2503.18565}]. This approach offers a compute-efficient path to developing Small Language Models (SLMs) that retain the sophisticated attention-like capabilities of larger models through recurrent mechanisms [\texttt{arXiv:2503.18565}].

\subsection{The Frontier: Interpretability, Design Principles, and Advanced Dynamics}

Recent research has shifted from merely scaling Transformer architectures toward a granular understanding of their internal mechanics, design requirements, and capacity for selective information processing. This frontier is characterized by efforts to disentangle the functional roles of attention and MLP layers, identify indispensable design principles, and enhance the model's ability to manage long-context dynamics through improved interpretability and structural modifications.

A significant portion of recent literature focuses on deconstructing the fundamental components of the attention mechanism. While standard models rely on complex mathematical forms, research suggests that certain elements are more critical than others; for instance, token mixing across positions is deemed indispensable, as its absence results in near-random behavior, whereas the specific mathematical form and sequence dependency can be substantially relaxed, particularly within hybrid architectures \texttt{arXiv:2510.11602}. This finding is complemented by investigations into the division of labor between layers, where attention is primarily responsible for in-context reasoning while MLPs handle knowledge storage \texttt{arXiv:2506.01115}. Furthermore, the role of attention in sequence modeling is challenged by the discovery that even frozen key and query weights can successfully form induction heads and perform competitive language modeling \texttt{arXiv:2506.01115}.

The ability of Transformers to navigate complex temporal structures is further elucidated through the study of induction heads. Research demonstrates that Transformers can dynamically select correct causal structures—specifically identifying the correct "lag" in interleaved Markov chains—to predict subsequent tokens \texttt{arXiv:2509.08184}. This capacity for selective focus is also being refined through training methodologies. For example, Adversarial Feedback for Attention (AFA) utilizes a policy gradient approach to redistribute attention from common words to task-relevant terms, yielding a 12.6\% performance improvement in large models \texttt{arXiv:2512.20661}. Similarly, "Self-Ablating Transformers" employ a k-winner-takes-all constraint during training to enforce selective activation, which increases neuron specialization and feature localization without compromising performance, though it notably decreases overall sparsity \texttt{arXiv:2505.00509}.

Finally, addressing the limitations of standard attention in long-context scenarios has led to structural innovations like the Forgetting Transformer (FoX). By integrating a data-dependent forget gate directly into the softmax mechanism via a logit bias, FoX bridges the gap between the high-retrieval capabilities of Transformers and the efficient forgetting seen in recurrent models \texttt{arXiv:2503.02130}. While FoX demonstrates superior length extrapolation and near-perfect accuracy in needle-in-the-haystack tests compared to models like Mamba-2 and DeltaNet, its scalability remains an open question, as empirical validation has been limited to models up to 760M parameters and 16k token contexts \texttt{arXiv:2503.02130}. These advancements collectively point toward a future of more efficient, interpretable, and contextually aware architectures \texttt{arXiv:2508.09834}.

\subsection{Tensions and Open Problems}

A primary tension in the literature concerns the fundamental necessity of the attention mechanism itself. While traditional paradigms emphasize the importance of specialized attention heads for capturing diverse linguistic features (\texttt{arXiv:1905.09418}), recent critical inquiries suggest that the heavy computational toll of attention may be misplaced. For instance, \texttt{arXiv:2506.11305} argues that the reliance on complex attention patterns—including those found in bidirectional models like BERT (\texttt{arXiv:1810.04805}) or rotary-enhanced architectures like RoFormer (\texttt{arXiv:2104.09864})—might be over-engineered, suggesting that much of the modeling power could be derived from simpler or more efficient structures. This debate is further complicated by the question of whether attention is even the primary driver of intelligence; as \texttt{arXiv:2505.06633} notes, the role of feedforward networks may be as critical as the attention mechanism, challenging the "attention-centric" view that has dominated the field for years.

Another significant disagreement exists regarding the optimal path toward long-context modeling and sequence efficiency. One school of thought advocates for approximating or sparsifying the attention matrix to maintain scalability, as seen in the linear complexity approach of Linformer (\texttt{arXiv:2006.04768}) or the adaptive span mechanisms of \texttt{arXiv:1905.07799}. However, these efficiency-driven models face scrutiny regarding their ability to maintain the representational density required for complex tasks. This tension is highlighted by the critique in \texttt{arXiv:2508.09834}, which weighs the speed gains of efficient architectures against the robust performance of standard dense models like BERT. Furthermore, while some researchers propose replacing attention with recurrent-like structures to achieve better efficiency (\texttt{arXiv:2006.16236}), others argue that such transitions may lose the specific synergistic computational graph effects that make modern Transformers so effective (\texttt{arXiv:2507.02944}).

Finally, there is a growing conflict between the pursuit of architectural complexity and the need for resource-efficient deployment. While large-scale pre-training on massive datasets has been the standard for maximizing performance (\texttt{arXiv:1806.00187}), recent work focuses on the "low-resource" frontier, where the goal is to adapt long-context capabilities to constrained environments. This has led to a divergence in methodology: some favor extending the context window through sparse or hierarchical attention, such as Longformer (\texttt{arXiv:2004.05150}), while others advocate for adaptive mechanisms that adjust attention spans dynamically (\texttt{arXiv:1905.07799}). As the field moves toward even larger contexts, the "curse of high dimensionality" (\texttt{arXiv:2505.22107}) presents a new open problem, questioning whether current attention-based scaling laws can hold as sequence lengths approach infinity.

\subsubsection{Outlook}

The frontier of language modeling is shifting from "more attention" to "smarter memory." As the industry moves beyond the simple scaling of parameters, the next generation of architectures will likely focus on the integration of persistent, compact memory structures (\texttt{arXiv:2505.00929}) and hybrid recurrent-attention mechanisms (\texttt{arXiv:2503.18565}) that can seamlessly transition between short-term reasoning and long-term context retention. The ultimate goal is to resolve the trade-off between computational speed and representational depth, moving toward models that do not just "pay attention," but intelligently decide what to remember and what to forget.

\section*{References}
\begin{enumerate}\small
  \item Learning Phrase Representations using RNN Encoder-Decoder for
  Statistical Machine Translation (2014). \texttt{arXiv:1406.1078}
  \item Neural Machine Translation by Jointly Learning to Align and Translate (2014). \texttt{arXiv:1409.0473}
  \item Sequence to Sequence Learning with Neural Networks (2014). \texttt{arXiv:1409.3215}
  \item Adam: A Method for Stochastic Optimization (2014). \texttt{arXiv:1412.6980}
  \item Neural Machine Translation of Rare Words with Subword Units (2015). \texttt{arXiv:1508.07909}
  \item Deep Residual Learning for Image Recognition (2015). \texttt{arXiv:1512.03385}
  \item Aligning Books and Movies: Towards Story-like Visual Explanations by
  Watching Movies and Reading Books (2015). \texttt{arXiv:1506.06724}
  \item Effective Approaches to Attention-based Neural Machine Translation (2015). \texttt{arXiv:1508.04025}
  \item Distilling the Knowledge in a Neural Network (2015). \texttt{arXiv:1503.02531}
  \item Layer Normalization (2016). \texttt{arXiv:1607.06450}
  \item SQuAD: 100,000+ Questions for Machine Comprehension of Text (2016). \texttt{arXiv:1606.05250}
  \item Google's Neural Machine Translation System: Bridging the Gap between
  Human and Machine Translation (2016). \texttt{arXiv:1609.08144}
  \item Pointer Sentinel Mixture Models (2016). \texttt{arXiv:1609.07843}
  \item Using the Output Embedding to Improve Language Models (2016). \texttt{arXiv:1608.05859}
  \item Can Active Memory Replace Attention? (2016). \texttt{arXiv:1610.08613}
  \item Attention Is All You Need (2017). \texttt{arXiv:1706.03762}
  \item Convolutional Sequence to Sequence Learning (2017). \texttt{arXiv:1705.03122}
  \item Outrageously Large Neural Networks: The Sparsely-Gated
  Mixture-of-Experts Layer (2017). \texttt{arXiv:1701.06538}
  \item A Broad-Coverage Challenge Corpus for Sentence Understanding through
  Inference (2017). \texttt{arXiv:1704.05426}
  \item Decoupled Weight Decay Regularization (2017). \texttt{arXiv:1711.05101}
  \item BERT: Pre-training of Deep Bidirectional Transformers for Language
  Understanding (2018). \texttt{arXiv:1810.04805}
  \item GLUE: A Multi-Task Benchmark and Analysis Platform for Natural Language
  Understanding (2018). \texttt{arXiv:1804.07461}
  \item Deep contextualized word representations (2018). \texttt{arXiv:1802.05365}
  \item Character-Level Language Modeling with Deeper Self-Attention (2018). \texttt{arXiv:1808.04444}
  \item Self-Attention with Relative Position Representations (2018). \texttt{arXiv:1803.02155}
  \item Adaptive Input Representations for Neural Language Modeling (2018). \texttt{arXiv:1809.10853}
  \item Scaling Neural Machine Translation (2018). \texttt{arXiv:1806.00187}
  \item Transformer-XL: Attentive Language Models Beyond a Fixed-Length Context (2019). \texttt{arXiv:1901.02860}
  \item RoBERTa: A Robustly Optimized BERT Pretraining Approach (2019). \texttt{arXiv:1907.11692}
  \item Generating Long Sequences with Sparse Transformers (2019). \texttt{arXiv:1904.10509}
  \item Exploring the Limits of Transfer Learning with a Unified Text-to-Text
  Transformer (2019). \texttt{arXiv:1910.10683}
  \item XLNet: Generalized Autoregressive Pretraining for Language Understanding (2019). \texttt{arXiv:1906.08237}
  \item ALBERT: A Lite BERT for Self-supervised Learning of Language
  Representations (2019). \texttt{arXiv:1909.11942}
  \item Analyzing Multi-Head Self-Attention: Specialized Heads Do the Heavy
  Lifting, the Rest Can Be Pruned (2019). \texttt{arXiv:1905.09418}
  \item HuggingFace's Transformers: State-of-the-art Natural Language Processing (2019). \texttt{arXiv:1910.03771}
  \item Adaptive Attention Span in Transformers (2019). \texttt{arXiv:1905.07799}
  \item Are Sixteen Heads Really Better than One? (2019). \texttt{arXiv:1905.10650}
  \item Compressive Transformers for Long-Range Sequence Modelling (2019). \texttt{arXiv:1911.05507}
  \item Language Models are Few-Shot Learners (2020). \texttt{arXiv:2005.14165}
  \item Longformer: The Long-Document Transformer (2020). \texttt{arXiv:2004.05150}
  \item Reformer: The Efficient Transformer (2020). \texttt{arXiv:2001.04451}
  \item An Image is Worth 16x16 Words: Transformers for Image Recognition at
  Scale (2020). \texttt{arXiv:2010.11929}
  \item Linformer: Self-Attention with Linear Complexity (2020). \texttt{arXiv:2006.04768}
  \item Transformers are RNNs: Fast Autoregressive Transformers with Linear
  Attention (2020). \texttt{arXiv:2006.16236}
  \item Big Bird: Transformers for Longer Sequences (2020). \texttt{arXiv:2007.14062}
  \item Rethinking Attention with Performers (2020). \texttt{arXiv:2009.14794}
  \item Scaling Laws for Neural Language Models (2020). \texttt{arXiv:2001.08361}
  \item RoFormer: Enhanced Transformer with Rotary Position Embedding (2021). \texttt{arXiv:2104.09864}
  \item Efficient Attention Mechanisms for Large Language Models: A Survey (2025). \texttt{arXiv:2507.19595}
  \item Forgetting Transformer: Softmax Attention with a Forget Gate (2025). \texttt{arXiv:2503.02130}
  \item Is Random Attention Sufficient for Sequence Modeling? Disentangling Trainable Components in the Transformer (2025). \texttt{arXiv:2506.01115}
  \item Don't Pay Attention (2025). \texttt{arXiv:2506.11305}
  \item Deconstructing Attention: Investigating Design Principles for Effective Language Modeling (2025). \texttt{arXiv:2510.11602}
  \item Speed Always Wins: A Survey on Efficient Architectures for Large Language Models (2025). \texttt{arXiv:2508.09834}
  \item Curse of High Dimensionality Issue in Transformer for Long-context Modeling (2025). \texttt{arXiv:2505.22107}
  \item Attention Is Not All You Need: The Importance of Feedforward Networks in Transformer Models (2025). \texttt{arXiv:2505.06633}
  \item InfLLM-V2: Dense-Sparse Switchable Attention for Seamless Short-to-Long Adaptation (2025). \texttt{arXiv:2509.24663}
  \item Distil-xLSTM: Learning Attention Mechanisms through Recurrent Structures (2025). \texttt{arXiv:2503.18565}
  \item Compact Recurrent Transformer with Persistent Memory (2025). \texttt{arXiv:2505.00929}
  \item SWAN-GPT: An Efficient and Scalable Approach for Long-Context Language
  Modeling (2025). \texttt{arXiv:2504.08719}
  \item Efficient Language Modeling for Low-Resource Settings with Hybrid
  RNN-Transformer Architectures (2025). \texttt{arXiv:2502.00617}
  \item StackTrans: From Large Language Model to Large Pushdown Automata Model (2025). \texttt{arXiv:2507.15343}
  \item Gated Associative Memory: A Parallel O(N) Architecture for Efficient Sequence Modeling (2025). \texttt{arXiv:2509.00605}
  \item PaPaformer: Language Model from Pre-trained Parallel Paths (2025). \texttt{arXiv:2508.00544}
  \item TPTT: Transforming Pretrained Transformers into Titans (2025). \texttt{arXiv:2506.17671}
  \item An extension of linear self-attention for in-context learning (2025). \texttt{arXiv:2503.23814}
  \item Inceptive Transformers: Enhancing Contextual Representations through Multi-Scale Feature Learning Across Domains and Languages (2025). \texttt{arXiv:2505.20496}
  \item Beyond Parallelism: Synergistic Computational Graph Effects in Multi-Head Attention (2025). \texttt{arXiv:2507.02944}
  \item Transformers as Multi-task Learners: Decoupling Features in Hidden Markov Models (2025). \texttt{arXiv:2506.01919}
  \item Parameter-Efficient Transformer Embeddings (2025). \texttt{arXiv:2505.02266}
  \item Short-Range Dependency Effects on Transformer Instability and a Decomposed Attention Solution (2025). \texttt{arXiv:2505.15548}
  \item Hierarchical Resolution Transformers: A Wavelet-Inspired Architecture for Multi-Scale Language Understanding (2025). \texttt{arXiv:2509.20581}
  \item Reservoir Computing as a Language Model (2025). \texttt{arXiv:2507.15779}
  \item Delayed Attention Training Improves Length Generalization in Transformer--RNN Hybrids (2025). \texttt{arXiv:2510.00258}
  \item Self-Ablating Transformers: More Interpretability, Less Sparsity (2025). \texttt{arXiv:2505.00509}
  \item Selective Induction Heads: How Transformers Select Causal Structures In Context (2025). \texttt{arXiv:2509.08184}
  \item ATLAS: Learning to Optimally Memorize the Context at Test Time (2025). \texttt{arXiv:2505.23735}
  \item Falcon-H1: A Family of Hybrid-Head Language Models Redefining Efficiency and Performance (2025). \texttt{arXiv:2507.22448}
  \item From Fake Focus to Real Precision: Confusion-Driven Adversarial Attention Learning in Transformers (2025). \texttt{arXiv:2512.20661}
  \item Mitigating Attention Localization in Small Scale: Self-Attention Refinement via One-step Belief Propagation (2025). \texttt{arXiv:2509.07324}
\end{enumerate}

\end{document}
